% =============================================================================
% Final LLM user prompt for per-sample rubric classification (Appendix)
% Requires \begin{promptbox}...\end{promptbox} defined in main preamble.
% Usage: % =============================================================================
% Final LLM user prompt for per-sample rubric classification (Appendix)
% Requires \begin{promptbox}...\end{promptbox} defined in main preamble.
% Usage: % =============================================================================
% Final LLM user prompt for per-sample rubric classification (Appendix)
% Requires \begin{promptbox}...\end{promptbox} defined in main preamble.
% Usage: % =============================================================================
% Final LLM user prompt for per-sample rubric classification (Appendix)
% Requires \begin{promptbox}...\end{promptbox} defined in main preamble.
% Usage: \input{appendix_rubric_prompt.tex}
% =============================================================================

\subsection{Per-Sample LLM Judge Prompt}
\label{app:rubric-prompt}

Each hold-out harmful sample triggers one API call with a single user message.
Surface-equal pairs ($s\!=\!t$) are pre-labeled as \texttt{consistent} by a deterministic
rule and are excluded from \texttt{\{\{ PAIR\_LINES \}\}}.
The string below is the exact \texttt{content} field sent to the judge
(\texttt{gpt-4o-mini}, temperature~$0$).
Only the four rubric blocks and the Input block are filled at runtime; Task and Output
are fixed.

\begin{promptbox}{LLM Judge Prompt for Token-Pair Classification}

## Task

You classify high-KL (student\_top1, teacher\_top1) token pairs for offline rule extraction.

We distill a student LLM toward a teacher that sees privileged safety guidance, while
the student sees only the user query. At high forward-KL positions in the student's
rollout, we compare the top-1 token preferred by the student versus the teacher.
Your job is to label each such pair with a semantic category so that we can aggregate
family-specific, query-universal lexicons (PIVOT\_STU/TEA, INTENT\_STU/TEA, RISK, FUNC,
and intent pairs) for training-time token filtering.

Use the user query and the student response below for context, but extract only tokens
that generalize across queries---not query-specific content nouns.

For each listed pair, assign **at most one** extractable category:
  pivot | intent | risk | function

If a pair does not belong to any of the four, use category=null.

Priority when ambiguous: pivot > risk > intent > function.

\vspace{1ex}

---
## Pivot

\texttt{\{\{ RUBRIC\_PIVOT \}\}}

\vspace{1ex}

---
## Intent

\texttt{\{\{ RUBRIC\_INTENT \}\}}

\vspace{1ex}

---
## Risk

\texttt{\{\{ RUBRIC\_RISK \}\}}

\vspace{1ex}

---
## Function

\texttt{\{\{ RUBRIC\_FUNCTION \}\}}

\vspace{1ex}

---
## Input

User query:

\texttt{\{\{ QUERY \}\}}

\vspace{1ex}

Student response:

\texttt{\{\{ STUDENT\_RESPONSE \}\}}

\vspace{1ex}

Top-\texttt{\{\{ N \}\}} high-KL token pairs (student\_top1 vs teacher\_top1):

\texttt{\{\{ PAIR\_LINES \}\}}

\vspace{1ex}

---
## Output

Return ONE JSON object (no markdown fence). Label exactly the pairs listed above.
Use rank to match input.

Extraction mapping:
  pivot   -> **pivot\_tokens**: only discourse/structural pivot cue(s) from this pair
            (student-side -> PIVOT\_STU, teacher-side -> PIVOT\_TEA)
  intent  -> student\_top1 -> INTENT\_STU, teacher\_top1 -> INTENT\_TEA;
            (student, teacher) -> \_INTENT\_PAIRS
            ONLY if stu->tea is a real safety-evaluation framing shift
            (not pronoun/person opener swaps)
  risk    -> **risk\_tokens**: only risk-related token(s) from this pair
            (one or both sides, but list risk token(s) only---not the whole pair)
  function-> **function\_tokens**: only function/stop token(s) from this pair

Schema:
\{
  "pair\_labels": [
    \{
      "rank": 1,
      "category": "pivot|intent|risk|function|null",
      "pivot\_tokens": [],
      "risk\_tokens": [],
      "function\_tokens": []
    \},
    ...
  ]
\}

Rules:
- One entry per listed rank only.
- category=null -> all token lists must be [].
- category=pivot -> pivot\_tokens must list only the pivot/transition cue(s) from the pair;
  omit non-pivot tokens.
- category=risk -> risk\_tokens must list only risk-related token(s) from the pair;
  omit non-risk tokens.
- category=function -> function\_tokens must be non-empty; list only function/stop
  token(s); omit content words.
- category=intent -> leave pivot\_tokens, risk\_tokens, and function\_tokens as [];
  label intent ONLY for stu->tea safety-evaluation framing (e.g., understand->recognize),
  NOT for The->I / Maybe->I / they->I style pronoun or person switches (use null).
- Unused token lists for a category must be [].

\end{promptbox}

\paragraph{Runtime bindings.}
\texttt{\{\{ RUBRIC\_PIVOT \}\}}, \texttt{\{\{ RUBRIC\_INTENT \}\}},
\texttt{\{\{ RUBRIC\_RISK \}\}}, and \texttt{\{\{ RUBRIC\_FUNCTION \}\}} are fixed
text (see the Pivot / Intent / Risk / Function rubric boxes in
Appendix~\ref{app:classifier:rubrics}).
\texttt{\{\{ QUERY \}\}} is the harmful user query;
\texttt{\{\{ STUDENT\_RESPONSE \}\}} is the student rollout;
\texttt{\{\{ N \}\}} is the number of disagreeing pairs shown;
\texttt{\{\{ PAIR\_LINES \}\}} lists one line per pair, e.g.:

\begin{promptbox}{Example pair line}
  rank=3  student\_top1='understand'  teacher\_top1='recognize'  forward\_kl=0.4521
\end{promptbox}

% =============================================================================

\subsection{Per-Sample LLM Judge Prompt}
\label{app:rubric-prompt}

Each hold-out harmful sample triggers one API call with a single user message.
Surface-equal pairs ($s\!=\!t$) are pre-labeled as \texttt{consistent} by a deterministic
rule and are excluded from \texttt{\{\{ PAIR\_LINES \}\}}.
The string below is the exact \texttt{content} field sent to the judge
(\texttt{gpt-4o-mini}, temperature~$0$).
Only the four rubric blocks and the Input block are filled at runtime; Task and Output
are fixed.

\begin{promptbox}{LLM Judge Prompt for Token-Pair Classification}

## Task

You classify high-KL (student\_top1, teacher\_top1) token pairs for offline rule extraction.

We distill a student LLM toward a teacher that sees privileged safety guidance, while
the student sees only the user query. At high forward-KL positions in the student's
rollout, we compare the top-1 token preferred by the student versus the teacher.
Your job is to label each such pair with a semantic category so that we can aggregate
family-specific, query-universal lexicons (PIVOT\_STU/TEA, INTENT\_STU/TEA, RISK, FUNC,
and intent pairs) for training-time token filtering.

Use the user query and the student response below for context, but extract only tokens
that generalize across queries---not query-specific content nouns.

For each listed pair, assign **at most one** extractable category:
  pivot | intent | risk | function

If a pair does not belong to any of the four, use category=null.

Priority when ambiguous: pivot > risk > intent > function.

\vspace{1ex}

---
## Pivot

\texttt{\{\{ RUBRIC\_PIVOT \}\}}

\vspace{1ex}

---
## Intent

\texttt{\{\{ RUBRIC\_INTENT \}\}}

\vspace{1ex}

---
## Risk

\texttt{\{\{ RUBRIC\_RISK \}\}}

\vspace{1ex}

---
## Function

\texttt{\{\{ RUBRIC\_FUNCTION \}\}}

\vspace{1ex}

---
## Input

User query:

\texttt{\{\{ QUERY \}\}}

\vspace{1ex}

Student response:

\texttt{\{\{ STUDENT\_RESPONSE \}\}}

\vspace{1ex}

Top-\texttt{\{\{ N \}\}} high-KL token pairs (student\_top1 vs teacher\_top1):

\texttt{\{\{ PAIR\_LINES \}\}}

\vspace{1ex}

---
## Output

Return ONE JSON object (no markdown fence). Label exactly the pairs listed above.
Use rank to match input.

Extraction mapping:
  pivot   -> **pivot\_tokens**: only discourse/structural pivot cue(s) from this pair
            (student-side -> PIVOT\_STU, teacher-side -> PIVOT\_TEA)
  intent  -> student\_top1 -> INTENT\_STU, teacher\_top1 -> INTENT\_TEA;
            (student, teacher) -> \_INTENT\_PAIRS
            ONLY if stu->tea is a real safety-evaluation framing shift
            (not pronoun/person opener swaps)
  risk    -> **risk\_tokens**: only risk-related token(s) from this pair
            (one or both sides, but list risk token(s) only---not the whole pair)
  function-> **function\_tokens**: only function/stop token(s) from this pair

Schema:
\{
  "pair\_labels": [
    \{
      "rank": 1,
      "category": "pivot|intent|risk|function|null",
      "pivot\_tokens": [],
      "risk\_tokens": [],
      "function\_tokens": []
    \},
    ...
  ]
\}

Rules:
- One entry per listed rank only.
- category=null -> all token lists must be [].
- category=pivot -> pivot\_tokens must list only the pivot/transition cue(s) from the pair;
  omit non-pivot tokens.
- category=risk -> risk\_tokens must list only risk-related token(s) from the pair;
  omit non-risk tokens.
- category=function -> function\_tokens must be non-empty; list only function/stop
  token(s); omit content words.
- category=intent -> leave pivot\_tokens, risk\_tokens, and function\_tokens as [];
  label intent ONLY for stu->tea safety-evaluation framing (e.g., understand->recognize),
  NOT for The->I / Maybe->I / they->I style pronoun or person switches (use null).
- Unused token lists for a category must be [].

\end{promptbox}

\paragraph{Runtime bindings.}
\texttt{\{\{ RUBRIC\_PIVOT \}\}}, \texttt{\{\{ RUBRIC\_INTENT \}\}},
\texttt{\{\{ RUBRIC\_RISK \}\}}, and \texttt{\{\{ RUBRIC\_FUNCTION \}\}} are fixed
text (see the Pivot / Intent / Risk / Function rubric boxes in
Appendix~\ref{app:classifier:rubrics}).
\texttt{\{\{ QUERY \}\}} is the harmful user query;
\texttt{\{\{ STUDENT\_RESPONSE \}\}} is the student rollout;
\texttt{\{\{ N \}\}} is the number of disagreeing pairs shown;
\texttt{\{\{ PAIR\_LINES \}\}} lists one line per pair, e.g.:

\begin{promptbox}{Example pair line}
  rank=3  student\_top1='understand'  teacher\_top1='recognize'  forward\_kl=0.4521
\end{promptbox}

% =============================================================================

\subsection{Per-Sample LLM Judge Prompt}
\label{app:rubric-prompt}

Each hold-out harmful sample triggers one API call with a single user message.
Surface-equal pairs ($s\!=\!t$) are pre-labeled as \texttt{consistent} by a deterministic
rule and are excluded from \texttt{\{\{ PAIR\_LINES \}\}}.
The string below is the exact \texttt{content} field sent to the judge
(\texttt{gpt-4o-mini}, temperature~$0$).
Only the four rubric blocks and the Input block are filled at runtime; Task and Output
are fixed.

\begin{promptbox}{LLM Judge Prompt for Token-Pair Classification}

## Task

You classify high-KL (student\_top1, teacher\_top1) token pairs for offline rule extraction.

We distill a student LLM toward a teacher that sees privileged safety guidance, while
the student sees only the user query. At high forward-KL positions in the student's
rollout, we compare the top-1 token preferred by the student versus the teacher.
Your job is to label each such pair with a semantic category so that we can aggregate
family-specific, query-universal lexicons (PIVOT\_STU/TEA, INTENT\_STU/TEA, RISK, FUNC,
and intent pairs) for training-time token filtering.

Use the user query and the student response below for context, but extract only tokens
that generalize across queries---not query-specific content nouns.

For each listed pair, assign **at most one** extractable category:
  pivot | intent | risk | function

If a pair does not belong to any of the four, use category=null.

Priority when ambiguous: pivot > risk > intent > function.

\vspace{1ex}

---
## Pivot

\texttt{\{\{ RUBRIC\_PIVOT \}\}}

\vspace{1ex}

---
## Intent

\texttt{\{\{ RUBRIC\_INTENT \}\}}

\vspace{1ex}

---
## Risk

\texttt{\{\{ RUBRIC\_RISK \}\}}

\vspace{1ex}

---
## Function

\texttt{\{\{ RUBRIC\_FUNCTION \}\}}

\vspace{1ex}

---
## Input

User query:

\texttt{\{\{ QUERY \}\}}

\vspace{1ex}

Student response:

\texttt{\{\{ STUDENT\_RESPONSE \}\}}

\vspace{1ex}

Top-\texttt{\{\{ N \}\}} high-KL token pairs (student\_top1 vs teacher\_top1):

\texttt{\{\{ PAIR\_LINES \}\}}

\vspace{1ex}

---
## Output

Return ONE JSON object (no markdown fence). Label exactly the pairs listed above.
Use rank to match input.

Extraction mapping:
  pivot   -> **pivot\_tokens**: only discourse/structural pivot cue(s) from this pair
            (student-side -> PIVOT\_STU, teacher-side -> PIVOT\_TEA)
  intent  -> student\_top1 -> INTENT\_STU, teacher\_top1 -> INTENT\_TEA;
            (student, teacher) -> \_INTENT\_PAIRS
            ONLY if stu->tea is a real safety-evaluation framing shift
            (not pronoun/person opener swaps)
  risk    -> **risk\_tokens**: only risk-related token(s) from this pair
            (one or both sides, but list risk token(s) only---not the whole pair)
  function-> **function\_tokens**: only function/stop token(s) from this pair

Schema:
\{
  "pair\_labels": [
    \{
      "rank": 1,
      "category": "pivot|intent|risk|function|null",
      "pivot\_tokens": [],
      "risk\_tokens": [],
      "function\_tokens": []
    \},
    ...
  ]
\}

Rules:
- One entry per listed rank only.
- category=null -> all token lists must be [].
- category=pivot -> pivot\_tokens must list only the pivot/transition cue(s) from the pair;
  omit non-pivot tokens.
- category=risk -> risk\_tokens must list only risk-related token(s) from the pair;
  omit non-risk tokens.
- category=function -> function\_tokens must be non-empty; list only function/stop
  token(s); omit content words.
- category=intent -> leave pivot\_tokens, risk\_tokens, and function\_tokens as [];
  label intent ONLY for stu->tea safety-evaluation framing (e.g., understand->recognize),
  NOT for The->I / Maybe->I / they->I style pronoun or person switches (use null).
- Unused token lists for a category must be [].

\end{promptbox}

\paragraph{Runtime bindings.}
\texttt{\{\{ RUBRIC\_PIVOT \}\}}, \texttt{\{\{ RUBRIC\_INTENT \}\}},
\texttt{\{\{ RUBRIC\_RISK \}\}}, and \texttt{\{\{ RUBRIC\_FUNCTION \}\}} are fixed
text (see the Pivot / Intent / Risk / Function rubric boxes in
Appendix~\ref{app:classifier:rubrics}).
\texttt{\{\{ QUERY \}\}} is the harmful user query;
\texttt{\{\{ STUDENT\_RESPONSE \}\}} is the student rollout;
\texttt{\{\{ N \}\}} is the number of disagreeing pairs shown;
\texttt{\{\{ PAIR\_LINES \}\}} lists one line per pair, e.g.:

\begin{promptbox}{Example pair line}
  rank=3  student\_top1='understand'  teacher\_top1='recognize'  forward\_kl=0.4521
\end{promptbox}

% =============================================================================

\subsection{Per-Sample LLM Judge Prompt}
\label{app:rubric-prompt}

Each hold-out harmful sample triggers one API call with a single user message.
Surface-equal pairs ($s\!=\!t$) are pre-labeled as \texttt{consistent} by a deterministic
rule and are excluded from \texttt{\{\{ PAIR\_LINES \}\}}.
The string below is the exact \texttt{content} field sent to the judge
(\texttt{gpt-4o-mini}, temperature~$0$).
Only the four rubric blocks and the Input block are filled at runtime; Task and Output
are fixed.

\begin{promptbox}{LLM Judge Prompt for Token-Pair Classification}

## Task

You classify high-KL (student\_top1, teacher\_top1) token pairs for offline rule extraction.

We distill a student LLM toward a teacher that sees privileged safety guidance, while
the student sees only the user query. At high forward-KL positions in the student's
rollout, we compare the top-1 token preferred by the student versus the teacher.
Your job is to label each such pair with a semantic category so that we can aggregate
family-specific, query-universal lexicons (PIVOT\_STU/TEA, INTENT\_STU/TEA, RISK, FUNC,
and intent pairs) for training-time token filtering.

Use the user query and the student response below for context, but extract only tokens
that generalize across queries---not query-specific content nouns.

For each listed pair, assign **at most one** extractable category:
  pivot | intent | risk | function

If a pair does not belong to any of the four, use category=null.

Priority when ambiguous: pivot > risk > intent > function.

\vspace{1ex}

---
## Pivot

\texttt{\{\{ RUBRIC\_PIVOT \}\}}

\vspace{1ex}

---
## Intent

\texttt{\{\{ RUBRIC\_INTENT \}\}}

\vspace{1ex}

---
## Risk

\texttt{\{\{ RUBRIC\_RISK \}\}}

\vspace{1ex}

---
## Function

\texttt{\{\{ RUBRIC\_FUNCTION \}\}}

\vspace{1ex}

---
## Input

User query:

\texttt{\{\{ QUERY \}\}}

\vspace{1ex}

Student response:

\texttt{\{\{ STUDENT\_RESPONSE \}\}}

\vspace{1ex}

Top-\texttt{\{\{ N \}\}} high-KL token pairs (student\_top1 vs teacher\_top1):

\texttt{\{\{ PAIR\_LINES \}\}}

\vspace{1ex}

---
## Output

Return ONE JSON object (no markdown fence). Label exactly the pairs listed above.
Use rank to match input.

Extraction mapping:
  pivot   -> **pivot\_tokens**: only discourse/structural pivot cue(s) from this pair
            (student-side -> PIVOT\_STU, teacher-side -> PIVOT\_TEA)
  intent  -> student\_top1 -> INTENT\_STU, teacher\_top1 -> INTENT\_TEA;
            (student, teacher) -> \_INTENT\_PAIRS
            ONLY if stu->tea is a real safety-evaluation framing shift
            (not pronoun/person opener swaps)
  risk    -> **risk\_tokens**: only risk-related token(s) from this pair
            (one or both sides, but list risk token(s) only---not the whole pair)
  function-> **function\_tokens**: only function/stop token(s) from this pair

Schema:
\{
  "pair\_labels": [
    \{
      "rank": 1,
      "category": "pivot|intent|risk|function|null",
      "pivot\_tokens": [],
      "risk\_tokens": [],
      "function\_tokens": []
    \},
    ...
  ]
\}

Rules:
- One entry per listed rank only.
- category=null -> all token lists must be [].
- category=pivot -> pivot\_tokens must list only the pivot/transition cue(s) from the pair;
  omit non-pivot tokens.
- category=risk -> risk\_tokens must list only risk-related token(s) from the pair;
  omit non-risk tokens.
- category=function -> function\_tokens must be non-empty; list only function/stop
  token(s); omit content words.
- category=intent -> leave pivot\_tokens, risk\_tokens, and function\_tokens as [];
  label intent ONLY for stu->tea safety-evaluation framing (e.g., understand->recognize),
  NOT for The->I / Maybe->I / they->I style pronoun or person switches (use null).
- Unused token lists for a category must be [].

\end{promptbox}

\paragraph{Runtime bindings.}
\texttt{\{\{ RUBRIC\_PIVOT \}\}}, \texttt{\{\{ RUBRIC\_INTENT \}\}},
\texttt{\{\{ RUBRIC\_RISK \}\}}, and \texttt{\{\{ RUBRIC\_FUNCTION \}\}} are fixed
text (see the Pivot / Intent / Risk / Function rubric boxes in
Appendix~\ref{app:classifier:rubrics}).
\texttt{\{\{ QUERY \}\}} is the harmful user query;
\texttt{\{\{ STUDENT\_RESPONSE \}\}} is the student rollout;
\texttt{\{\{ N \}\}} is the number of disagreeing pairs shown;
\texttt{\{\{ PAIR\_LINES \}\}} lists one line per pair, e.g.:

\begin{promptbox}{Example pair line}
  rank=3  student\_top1='understand'  teacher\_top1='recognize'  forward\_kl=0.4521
\end{promptbox}
